Investiga y compara alternativas de sistemas operativos a Kali Linux. ¿Cuál es su principal propósito?

Algunas alternativas de Kali Linux en el área de ciberseguridad son Parrot Security OS (un sistema con un enfoque más amigable y de uso diario que Kali), BlackArch Linux (una distribución basada en Arch con un repositorio masivo de herramientas) y BackBox Linux (una distribución basada en Ubuntu ligera y optimizada para pentesting).

Recopilé algunas de sus características en el siguiente cuadro comparativo:

\renewcommand{\arraystretch}{1.3}
\begin{table}[H]
    \centering
    \footnotesize
    \begin{tabularx}{\textwidth}{|>{\bfseries\raggedright\arraybackslash}p{2.8cm}|Y|Y|Y|Y|}
    \hline
    \textbf{Característica} & \textbf{Kali Linux} & \textbf{Parrot Security OS} & \textbf{BlackArch Linux} & \textbf{BackBox Linux} \\
    \hline
    Basado en & Debian Testing & Debian & Arch Linux & Ubuntu \\
    \hline
    Enfoque Principal & Pentesting puro y análisis forense. & Pentesting, privacidad y desarrollo. Uso diario. & Pentesting avanzado e investigación de seguridad. & Auditorías de seguridad y pentesting. \\
    \hline
    Público Objetivo & Profesionales de seguridad y estudiantes intermedios. & Principiantes y profesionales (muy amigable). & Usuarios muy avanzados y expertos en seguridad. & Profesionales de seguridad y administradores. \\
    \hline
    Entorno de Escritorio & XFCE (soporta KDE, GNOME) & MATE (soporta KDE) & Gestores de ventanas (Fluxbox, Openbox, etc.) & XFCE \\
    \hline
    Herramientas & +600 preinstaladas & Cientos (incluye herramientas de anonimato/cripto) & +2,800 preinstaladas (repositorio masivo) & $\sim$70 (solo las más esenciales) \\
    \hline
    Consumo de Recursos & Moderado & Muy ligero / Eficiente & Variable (ligero en RAM, ISO completa pesada) & Muy ligero (ideal para hardware modesto) \\
    \hline
    Gestor de Paquetes & APT (\texttt{apt}) & APT (\texttt{apt}) & Pacman (\texttt{pacman}) & APT (\texttt{apt}) \\
    \end{tabularx}
\end{table}

En general, encontré que Parrot Security OS es la alternativa más reconocida dentro del área por ser una distribución más amigable que Kali; además de venir con herramientas de hacking de manera nativa, el sistema viene configurado con herramientas para proteger la privacidad como el enrutamiento a través de Tor, y también es más recomendable para el uso diario.
Por el contrario, BlackArch tiene una gran cantidad de herramientas para realizar tareas de seguridad y, al estar basado en Arch Linux, es un sistema \textit{rolling release} puro y se compila de manera más personalizada.
Finalmente, BackBox Linux es una opción más ligera, estable y rápida, pues cuenta con menos herramientas (aquellas que probablemente no utilicen usuarios menos experimentados), sino que incluye únicamente las herramientas más reconocidas, útiles y probadas para garantizar la velocidad y la estabilidad.
